\documentclass[a4paper,11pt]{amsart}

\usepackage[margin=1in]{geometry}
\usepackage{amsmath,amssymb,amsthm,mathtools}
\usepackage{enumitem}
\usepackage[colorlinks=true,linkcolor=blue,citecolor=blue,urlcolor=blue]{hyperref}

\newtheorem{theorem}{Theorem}[section]
\newtheorem{proposition}[theorem]{Proposition}
\newtheorem{lemma}[theorem]{Lemma}
\newtheorem{corollary}[theorem]{Corollary}
\newtheorem{importedtheorem}[theorem]{Imported theorem}
\theoremstyle{remark}
\newtheorem{remark}[theorem]{Remark}

\newcommand{\Sub}{\operatorname{Sub}}
\newcommand{\Sym}{\operatorname{Sym}}
\newcommand{\F}{\mathbb F}
\newcommand{\OO}{\mathcal O}
\newcommand{\Prob}{\mathbb P}
\newcommand{\E}{\mathbb E}
\newcommand{\gbinom}[2]{\genfrac{[}{]}{0pt}{}{#1}{#2}_2}

\title[Subgroups of the symmetric group]{The logarithmic number of subgroups of the symmetric group}
\author{}
\date{}

\begin{document}

\begin{abstract}
Let
\[
        s(n)=|\{H:H\leq S_n\}|.
\]
The main asymptotic formula is
\[
        \log_2 s(n)=\frac{n^2}{16}+o(n^2),
        \qquad\text{equivalently}\qquad
        s(n)=2^{(1/16+o(1))n^2}.
\]
The lower bound is elementary and comes from the subgroup lattice of an elementary abelian
2-subgroup of rank \(\lfloor n/2\rfloor\). The matching upper bound is the deep theorem of
Pyber on subgroup enumeration in symmetric groups, and is explicitly quoted as an imported
theorem rather than reproved here. We also prove the exact statistical statement in the elementary
abelian model: a uniformly random subgroup of \(C_2^{\lfloor n/2\rfloor}\leq S_n\) has order
\(2^{n/4+o(n)}\). Finally, we record a caution: the leading asymptotic alone does not determine
the order of a uniformly random subgroup of \(S_n\); a dihedral-block construction gives
\(2^{n^2/16+O(n)}\) subgroups of order \(2^{n/2+o(n)}\) when \(4\mid n\).
\end{abstract}

\maketitle

\section{Statement of the problem and scope}

All logarithms are base \(2\). Let
\[
        s(n)=|\Sub(S_n)|
\]
denote the number of subgroups of the symmetric group on \(n\) letters. The principal theorem is
as follows.

\begin{theorem}\label{thm:main}
As \(n\to\infty\),
\[
        s(n)=2^{n^2/16+o(n^2)}.
\]
Equivalently,
\[
        \log_2 s(n)\sim \frac{n^2}{16}.
\]
\end{theorem}

The proof has two logically different parts. The lower bound is elementary and is proved in full
below. The upper bound is a substantial finite-permutation-group theorem of Pyber. Since the
proof of Pyber's theorem is long and uses nontrivial subgroup-enumeration machinery, it is stated
in this note as an imported theorem. Thus every local lemma and proposition below is proved from
first principles, while the global upper bound is quoted precisely as the single external input.

The statistical question is more delicate. The elementary lower-bound model has a clean and fully
proved order law: a random subspace of \(\F_2^{\lfloor n/2\rfloor}\) has dimension
\(n/4+o(n)\), hence order \(2^{n/4+o(n)}\). However, this model theorem should not be confused
with a theorem about a uniformly random subgroup of all of \(S_n\). Section~\ref{sec:dihedral}
shows that there are also \(2^{n^2/16+O(n)}\) subgroups of \(S_n\) with order
\(2^{n/2+o(n)}\) for \(4\mid n\). Consequently, the leading logarithmic formula by itself is too
coarse to determine the global order distribution.

\section{Gaussian binomial coefficients}

For integers \(0\leq d\leq r\), let
\[
        \gbinom{r}{d}
\]
denote the number of \(d\)-dimensional linear subspaces of \(\F_2^r\). These numbers are the
Gaussian binomial coefficients over \(\F_2\).

\begin{lemma}[Formula for \(\gbinom{r}{d}\)]\label{lem:gauss-formula}
For \(0\leq d\leq r\),
\[
        \gbinom{r}{d}
        =
        \frac{(2^r-1)(2^r-2)\cdots(2^r-2^{d-1})}
        {(2^d-1)(2^d-2)\cdots(2^d-2^{d-1})}.
\]
\end{lemma}

\begin{proof}
Count ordered bases of \(d\)-dimensional subspaces of \(\F_2^r\) in two ways. An ordered
linearly independent \(d\)-tuple in \(\F_2^r\) can be chosen in
\[
        (2^r-1)(2^r-2)\cdots(2^r-2^{d-1})
\]
ways: after \(i\) independent vectors have been chosen, their span has \(2^i\) elements, so the
next vector has \(2^r-2^i\) choices. On the other hand, each fixed \(d\)-dimensional subspace has
\[
        (2^d-1)(2^d-2)\cdots(2^d-2^{d-1})
\]
ordered bases, by the same argument inside \(\F_2^d\). Dividing gives the formula.
\end{proof}

\begin{lemma}[Uniform Gaussian-binomial estimate]\label{lem:gauss-estimate}
There is an absolute constant
\[
        C_2=\prod_{j=1}^{\infty}(1-2^{-j})^{-1}<\infty
\]
such that, for all \(0\leq d\leq r\),
\[
        2^{d(r-d)}\leq \gbinom{r}{d}\leq C_2\,2^{d(r-d)}.
\]
\end{lemma}

\begin{proof}
Starting from Lemma~\ref{lem:gauss-formula}, factor out the powers of \(2\):
\[
        \gbinom{r}{d}
        =2^{d(r-d)}
        \prod_{i=0}^{d-1}\frac{1-2^{i-r}}{1-2^{i-d}}.
\]
Since \(r\geq d\), one has \(2^{i-r}\leq 2^{i-d}\), hence
\[
        1-2^{i-r}\geq 1-2^{i-d}.
\]
Each factor in the product is therefore at least \(1\), proving the lower bound.

For the upper bound, use \(1-2^{i-r}\leq 1\). Then
\[
        \prod_{i=0}^{d-1}\frac{1-2^{i-r}}{1-2^{i-d}}
        \leq
        \prod_{i=0}^{d-1}(1-2^{i-d})^{-1}
        =
        \prod_{j=1}^{d}(1-2^{-j})^{-1}
        \leq
        \prod_{j=1}^{\infty}(1-2^{-j})^{-1}.
\]
The infinite product converges because \(\sum_{j\geq 1}2^{-j}<\infty\). This proves the estimate.
\end{proof}

\begin{lemma}[Number of subspaces of \(\F_2^r\)]\label{lem:subspace-total}
Let
\[
        G(r)=\sum_{d=0}^{r}\gbinom{r}{d}.
\]
Then
\[
        G(r)=2^{r^2/4+O(\log r)}.
\]
In particular,
\[
        \log_2 G(r)=\frac{r^2}{4}+O(\log r).
\]
\end{lemma}

\begin{proof}
The function \(d(r-d)\) is maximized at \(d=r/2\), with maximum \(r^2/4\). Therefore
Lemma~\ref{lem:gauss-estimate} gives
\[
        G(r)
        \leq
        \sum_{d=0}^r C_2 2^{d(r-d)}
        \leq
        (r+1)C_2 2^{r^2/4}.
\]
For the lower bound, take \(d=\lfloor r/2\rfloor\). Then
\[
        G(r)\geq \gbinom{r}{\lfloor r/2\rfloor}
        \geq 2^{\lfloor r/2\rfloor\lceil r/2\rceil}
        =2^{r^2/4+O(1)}.
\]
Combining the two estimates gives the claim.
\end{proof}

\begin{lemma}[Concentration of the dimension of a random subspace]\label{lem:dimension-concentration}
Let \(U_r\) be chosen uniformly from all linear subspaces of \(\F_2^r\). Then
\[
        \frac{\dim U_r}{r}\to \frac12
\]
in probability. Equivalently, for every fixed \(\varepsilon>0\),
\[
        \Prob\left(\left|\dim U_r-\frac r2\right|\geq \varepsilon r\right)\to 0.
\]
More quantitatively,
\[
        \Prob\left(\left|\dim U_r-\frac r2\right|\geq \varepsilon r\right)
        \leq
        (r+1)C_2\,2^{-\varepsilon^2r^2+O(1)}.
\]
\end{lemma}

\begin{proof}
For every \(d\),
\[
        d(r-d)=\frac{r^2}{4}-\left(d-\frac r2\right)^2.
\]
Thus, if \(|d-r/2|\geq \varepsilon r\), Lemma~\ref{lem:gauss-estimate} gives
\[
        \gbinom{r}{d}\leq C_2 2^{r^2/4-\varepsilon^2r^2}.
\]
There are at most \(r+1\) possible dimensions \(d\), so the number of exceptional subspaces is at
most
\[
        (r+1)C_2 2^{r^2/4-\varepsilon^2r^2}.
\]
By Lemma~\ref{lem:subspace-total}, the total number of subspaces is at least
\(2^{r^2/4+O(1)}\). Dividing the two bounds proves the quantitative estimate and hence the
probabilistic convergence.
\end{proof}

\section{The elementary lower bound for \texorpdfstring{\(s(n)\)}{s(n)}}

Let
\[
        m=\left\lfloor \frac n2\right\rfloor.
\]
Inside \(S_n\), define
\[
        E_m=\langle (1\,2),(3\,4),\ldots,(2m-1\,\,2m)\rangle.
\]
Then \(E_m\cong C_2^m\), because the displayed transpositions are disjoint and independent.

\begin{proposition}[Elementary lower bound]\label{prop:lower}
One has
\[
        s(n)\geq 2^{n^2/16+O(n)}.
\]
In fact,
\[
        s(n)\geq G\left(\left\lfloor \frac n2\right\rfloor\right),
\]
where \(G(r)\) is the number of subspaces of \(\F_2^r\).
\end{proposition}

\begin{proof}
The subgroup \(E_m\) is naturally an \(m\)-dimensional vector space over \(\F_2\): multiplication
in \(E_m\) corresponds to addition of exponent vectors in \(\F_2^m\). Therefore the subgroups of
\(E_m\) are exactly the \(\F_2\)-linear subspaces of \(\F_2^m\). Distinct subspaces give distinct
subgroups of \(S_n\). Hence
\[
        s(n)\geq |\Sub(E_m)|=G(m).
\]
By Lemma~\ref{lem:subspace-total},
\[
        G(m)=2^{m^2/4+O(\log m)}.
\]
Since \(m=\lfloor n/2\rfloor\), we have
\[
        \frac{m^2}{4}=\frac{n^2}{16}+O(n).
\]
Thus
\[
        s(n)\geq 2^{n^2/16+O(n)}.
\]
\end{proof}

\section{The upper bound and the asymptotic formula}

The following theorem is the non-elementary input.

\begin{importedtheorem}[Pyber's upper bound]\label{thm:pyber-upper}
As \(n\to\infty\),
\[
        s(n)\leq 2^{n^2/16+o(n^2)}.
\]
\end{importedtheorem}

\begin{remark}
Theorem~\ref{thm:pyber-upper} is the difficult half of the subgroup-enumeration theorem for
symmetric groups. It is not an elementary consequence of the lower-bound construction above. Its
proof uses substantial information about the structure and enumeration of finite permutation
groups. In this manuscript it is quoted as a black box; the purpose of the surrounding arguments is
to make the reduction to this single external theorem completely explicit.
\end{remark}

\begin{proof}[Proof of Theorem~\ref{thm:main}]
The lower bound from Proposition~\ref{prop:lower} gives
\[
        s(n)\geq 2^{n^2/16+O(n)}=2^{n^2/16+o(n^2)}.
\]
The imported upper bound, Theorem~\ref{thm:pyber-upper}, gives
\[
        s(n)\leq 2^{n^2/16+o(n^2)}.
\]
Combining the two inequalities yields
\[
        s(n)=2^{n^2/16+o(n^2)}.
\]
Taking logarithms gives
\[
        \log_2 s(n)=\frac{n^2}{16}+o(n^2),
\]
as claimed.
\end{proof}

\section{The order law inside the elementary abelian model}\label{sec:model-order}

The lower-bound construction not only gives many subgroups; it also gives a clean order statistic
inside that model.

\begin{theorem}[Model order theorem]\label{thm:model-order}
Let \(m=\lfloor n/2\rfloor\), and choose a subgroup \(U_n\leq E_m\) uniformly at random from all
subgroups of \(E_m\). Then
\[
        \log_2 |U_n|=\frac n4+o(n)
\]
in probability. Equivalently,
\[
        |U_n|=2^{n/4+o(n)}
\]
in probability.
\end{theorem}

\begin{proof}
Under the identification \(E_m\cong \F_2^m\), the subgroup \(U_n\) is a uniformly random subspace
of \(\F_2^m\). Its group order is
\[
        |U_n|=2^{\dim U_n},
\]
so
\[
        \log_2 |U_n|=\dim U_n.
\]
By Lemma~\ref{lem:dimension-concentration},
\[
        \dim U_n=\frac m2+o(m)
\]
in probability. Since \(m=n/2+O(1)\), this becomes
\[
        \log_2 |U_n|=\frac n4+o(n),
\]
which is the claim.
\end{proof}

\section{Why the global order distribution is a finer question}\label{sec:dihedral}

It is tempting to infer from Theorem~\ref{thm:model-order} that a uniformly random subgroup of
\(S_n\) itself should have order \(2^{n/4+o(n)}\). That inference is not justified by the leading
asymptotic formula. The following construction shows why: another natural family has the same
leading \(2^{n^2/16}\)-scale but a different linear order scale.

\begin{lemma}[A dihedral block group]\label{lem:dihedral-block}
Let \(D_8\) be the dihedral group of order \(8\), acting faithfully on four points as the symmetry
group of a square. Its Frattini subgroup \(\Phi(D_8)\), namely the intersection of all maximal
subgroups of \(D_8\), has order \(2\), and
\[
        D_8/\Phi(D_8)\cong C_2^2.
\]
\end{lemma}

\begin{proof}
Write
\[
        D_8=\langle r,s:r^4=s^2=1,\;srs=r^{-1}\rangle.
\]
The subgroup \(\langle r^2\rangle\) is central of order \(2\). The quotient
\(D_8/\langle r^2\rangle\) is generated by the images of \(r\) and \(s\), both of order \(2\), and
these images commute because the commutator relation becomes trivial modulo \(\langle r^2\rangle\).
Thus the quotient has order \(4\) and is isomorphic to \(C_2^2\).

It remains to identify \(\langle r^2\rangle\) with the Frattini subgroup. Since \(|D_8|=8\), the
maximal subgroups have order \(4\). They are exactly
\[
        \langle r\rangle,
        \qquad
        \langle r^2,s\rangle,
        \qquad
        \langle r^2,rs\rangle.
\]
Indeed, \(\langle r\rangle\) is the unique cyclic subgroup of order \(4\), while every subgroup of
order \(4\) not equal to \(\langle r\rangle\) contains a reflection and \(r^2\), giving one of the
two displayed Klein four subgroups according to whether the reflection is \(r^{2j}s\) or
\(r^{2j+1}s\). Their intersection is precisely \(\langle r^2\rangle\). Hence
\[
        \Phi(D_8)=\langle r\rangle\cap \langle r^2,s\rangle\cap\langle r^2,rs\rangle
        =\langle r^2\rangle.
\]
The quotient by \(\Phi(D_8)\) is therefore \(C_2^2\), as proved above.
\end{proof}

\begin{proposition}[Many subgroups at order scale \(2^{n/2}\)]\label{prop:dihedral-family}
Suppose \(n=4k\). Then \(S_n\) contains a family \(\mathcal F_n\) of subgroups with
\[
        |\mathcal F_n|=2^{n^2/16+O(\log n)}
\]
such that all but an exponentially small proportion of the members of \(\mathcal F_n\) satisfy
\[
        \log_2 |H|=\frac n2+o(n).
\]
\end{proposition}

\begin{proof}
Partition \(\{1,\ldots,n\}\) into \(k\) fixed blocks of size \(4\). On each block let a copy of
\(D_8\) act as in Lemma~\ref{lem:dihedral-block}, and set
\[
        P_k=D_8^k\leq S_{4k}.
\]
Let
\[
        N_k=\Phi(D_8)^k\trianglelefteq P_k.
\]
Then
\[
        N_k\cong C_2^k
        \qquad\text{and}\qquad
        P_k/N_k\cong C_2^{2k}.
\]
Let
\[
        \pi:P_k\to P_k/N_k
\]
be the quotient map. For every subspace \(U\leq P_k/N_k\), the inverse image
\[
        H_U=\pi^{-1}(U)
\]
is a subgroup of \(P_k\), hence a subgroup of \(S_{4k}\). Distinct subspaces \(U\) give distinct
subgroups \(H_U\), since \(\pi(H_U)=U\). Therefore the number of such subgroups is exactly the
number of subspaces of \(\F_2^{2k}\), namely
\[
        G(2k)=2^{(2k)^2/4+O(\log k)}=2^{k^2+O(\log k)}.
\]
Since \(n=4k\), this is
\[
        2^{n^2/16+O(\log n)}.
\]

Moreover, if \(\dim U=d\), then
\[
        |H_U|=|N_k|\,|U|=2^k 2^d=2^{k+d}.
\]
By Lemma~\ref{lem:dimension-concentration}, a uniformly random subspace \(U\leq\F_2^{2k}\) has
\[
        d=k+o(k)
\]
in probability. Hence, for all but an exponentially small proportion of the subspaces \(U\),
\[
        \log_2 |H_U|=k+d=2k+o(k)=\frac n2+o(n).
\]
This proves the claim.
\end{proof}

\begin{corollary}[The leading formula alone does not give a global order law]\label{cor:order-caution}
The asymptotic formula
\[
        s(n)=2^{n^2/16+o(n^2)}
\]
does not by itself imply that a uniformly random subgroup of \(S_n\) has order
\(2^{n/4+o(n)}\). Indeed, for \(4\mid n\), there is a family of
\[
        2^{n^2/16+O(\log n)}
\]
subgroups of order scale \(2^{n/2+o(n)}\), as shown in
Proposition~\ref{prop:dihedral-family}.
\end{corollary}

\begin{proof}
Theorem~\ref{thm:model-order} proves the order scale \(2^{n/4+o(n)}\) only for subgroups chosen
inside a fixed elementary abelian subgroup of rank \(\lfloor n/2\rfloor\). Proposition~\ref{prop:dihedral-family}
constructs, on the same leading \(2^{n^2/16}\) logarithmic scale, subgroups whose order is instead
\(2^{n/2+o(n)}\). Therefore the leading exponential count cannot distinguish the global
probability weights of these different families. A genuine theorem for a uniformly random subgroup
of \(S_n\) requires finer enumeration than the main logarithmic asymptotic.
\end{proof}

\section{A useful rank check for elementary abelian 2-subgroups}

The following elementary observation explains why the rank \(\lfloor n/2\rfloor\) appearing in the
lower bound is the correct maximum possible rank for elementary abelian \(2\)-subgroups of \(S_n\).

\begin{lemma}\label{lem:rank-bound}
If \(A\leq S_n\) is elementary abelian of exponent \(2\), say \(A\cong C_2^r\), then
\[
        r\leq \left\lfloor\frac n2\right\rfloor.
\]
\end{lemma}

\begin{proof}
Let \(A\) act on \(\{1,
\ldots,n\}\), and decompose the set into \(A\)-orbits
\[
        \Omega_1,\ldots,\Omega_t.
\]
For each orbit \(\Omega_i\), the image of \(A\) in \(\Sym(\Omega_i)\) is a transitive abelian
permutation group. A transitive abelian permutation group acts regularly after factoring out its
kernel: if an element fixes one point of a transitive orbit, then, because the group is abelian, it
fixes every point of that orbit. Hence the faithful image on \(\Omega_i\) is an elementary abelian
2-group of order \(|\Omega_i|\). Thus \(|\Omega_i|=2^{a_i}\) for some \(a_i\geq 0\), and the rank of
the image on \(\Omega_i\) is \(a_i\).

The diagonal action of \(A\) on all its orbits is faithful, so the natural homomorphism from \(A\)
to the product of its images on the orbits is injective. Therefore
\[
        r\leq \sum_{i=1}^t a_i.
\]
For each orbit, \(a_i\leq 2^{a_i}/2\) when \(a_i\geq 1\), and the inequality is also true for
\(a_i=0\). Hence
\[
        r\leq \sum_{i=1}^t a_i
        \leq
        \frac12\sum_{i=1}^t 2^{a_i}
        =\frac n2.
\]
Since \(r\) is an integer, \(r\leq\lfloor n/2\rfloor\).
\end{proof}

\section{Verification of the proof}

We finish by isolating exactly what has been proved and where each input enters.

\begin{enumerate}[label=\textup{(\arabic*)},leftmargin=2.5em]
    \item Lemmas~\ref{lem:gauss-formula}--\ref{lem:dimension-concentration} are fully elementary
    linear algebra over \(\F_2\). They prove both the count and dimension concentration for
    subspaces of \(\F_2^r\).

    \item Proposition~\ref{prop:lower} injects the subgroup lattice of
    \(C_2^{\lfloor n/2\rfloor}\) into \(\Sub(S_n)\). This proves the lower bound
    \[
            s(n)\geq 2^{n^2/16+O(n)}.
    \]

    \item The only imported result is Theorem~\ref{thm:pyber-upper}, Pyber's upper bound
    \[
            s(n)\leq 2^{n^2/16+o(n^2)}.
    \]
    Combining it with the elementary lower bound proves
    \[
            s(n)=2^{n^2/16+o(n^2)}.
    \]

    \item Theorem~\ref{thm:model-order} is a fully proved statistical theorem for the elementary
    abelian lower-bound model. It says that a random subgroup of
    \(C_2^{\lfloor n/2\rfloor}\leq S_n\) has order \(2^{n/4+o(n)}\).

    \item Proposition~\ref{prop:dihedral-family} proves that the global order distribution of a
    uniformly random subgroup of \(S_n\) cannot be read off from the leading logarithmic formula
    alone: for \(4\mid n\), there is also a family of \(2^{n^2/16+O(\log n)}\) subgroups with order
    \(2^{n/2+o(n)}\).
\end{enumerate}

Thus the asymptotic formula for the number of subgroups is rigorous once Pyber's upper-bound
theorem is admitted, while the order-distribution question is explicitly separated from the
coarser logarithmic enumeration.

\begin{thebibliography}{9}

\bibitem{Pyber}
L. Pyber,
\emph{The number of subgroups of the symmetric group},
Combinatorica \textbf{13} (1993), 363--372.

\end{thebibliography}

\end{document}
